\section{The complete invariant quadratic arithmetic target}
\label{sec:fr-Y2-target}

Put
\[
 H=G_{\mathbf Q},\qquad
 K_2=\mathbf Q(\mu_\infty,2^{1/\infty}),\qquad
 K=G_{K_2},\qquad
 Y_2=H^1_{\mathrm{cts}}(K,\mathbf F_2)^H.
\]
All cohomology below is continuous, with trivial coefficient module
\(\mathbf F_2\), unless otherwise indicated.
We use the established description
\[
 A=H/K
 =\{(u,t)\in U\ltimes\widehat{\mathbf Z}:
                  t\bmod2=\epsilon_8(u)\},
 \qquad U=\widehat{\mathbf Z}^{\times},
\]
with multiplication
\((u,t)(v,s)=(uv,t+us)\).
Every quadratic character of \(H\) factors through \(U\).
We also use the already established arithmetic characters
\[
 \psi_p=\rho_p^I/2\bmod2\quad(p\text{ odd prime}),\qquad
 \eta_3=m_3/8\bmod2
\]
on \(K\), their \(H\)-invariance, and the cubic identities recalled
in the proof below.

For an odd prime \(p\), set
\[
 p^*=(-1)^{(p-1)/2}p,\qquad
 \varepsilon_p=\frac{1-(2/p)}2\in\mathbf F_2,
 \qquad
 \ell_p(d)=\frac{1-(d/p)}2\in\mathbf F_2
 \quad(p\nmid d).
\]
Here \((d/p)\) is the Legendre symbol.
Let \(b,a,c_p\in H^1(U,\mathbf F_2)\) be the quadratic
characters of \(\mathbf Q(i)\), \(\mathbf Q(\sqrt2)\),
and \(\mathbf Q(\sqrt{p^*})\), respectively.
Thus
\[
 b(u)=\frac{u_2-1}{2}\bmod2,\qquad
 a(u)=\frac{u_2^2-1}{8}\bmod2.
\]

For each odd prime \(p\), let
\(k_p=v_2(p-1)\) and let \(\delta_p\in H^2(U,\mathbf F_2)\)
be the pullback of the unique nonzero class of
\[
 H^2(C_{2^{k_p}},\mathbf F_2)
\]
from the maximal \(2\)-primary cyclic quotient of
\(\mathbf Z_p^\times\).
Equivalently, it is the class of
\(C_{2^{k_p+1}}\to C_{2^{k_p}}\) with central kernel \(C_2\).
This class does not depend on a choice of cyclic generator.
Put \(\delta_2=b^2\).

\begin{theorem}\label{thm:fr-Y2-complete}
There is a direct-sum decomposition
\begin{equation}\label{eq:fr-Y2-basis}
 Y_2=
 \mathbf F_2\eta_3
 \ \oplus\!
 \bigoplus_{p\ {\rm odd}}\mathbf F_2\psi_p
 \ \oplus\!
 \bigoplus_{p\ {\rm odd}}\mathbf F_2\kappa_{2,p}
 \ \oplus\!
 \bigoplus_{\substack{p<q\\p,q\ {\rm odd}}}
                     \mathbf F_2\kappa_{p,q}.
\end{equation}
The additional characters are specified uniquely as follows.
The transgression gives an isomorphism
\begin{equation}\label{eq:fr-Y2-transgression}
 \operatorname{tr}:Y_2
 \xrightarrow{\ \sim\ }
 \ker\!\left(H^2(A,\mathbf F_2)
                   \longrightarrow\operatorname{Br}(\mathbf Q)[2]\right).
\end{equation}
Writing \(\operatorname{Inf}_U^A\) for inflation, one has
\begin{align}
 \operatorname{tr}(\kappa_{2,p})
 &=\operatorname{Inf}_U^A
       \left(a c_p+\varepsilon_p(\delta_2+\delta_p)\right),
                                                        \label{eq:fr-Y2-k2p}\\
 \operatorname{tr}(\kappa_{p,q})
 &=\operatorname{Inf}_U^A
       \left(c_pc_q+\ell_p(q^*)\delta_p
                       +\ell_q(p^*)\delta_q\right).
                                                        \label{eq:fr-Y2-kpq}
\end{align}
All products in these formulas are cup products.
Every sum in \eqref{eq:fr-Y2-basis} is finite.

In particular the quotient of \(Y_2\) by the known span
\(\mathbf F_2\eta_3+\bigoplus_p\mathbf F_2\psi_p\)
has countably infinite dimension.
These are statements about the arithmetic target only.
They do not assert that any additional basis character
extends to an invariant character of the corresponding full-GT subgroup.
\end{theorem}

\subsection{Transgression and the affine quotient}

The five-term sequence for \(1\to K\to H\to A\to1\) begins
\[
 0\longrightarrow H^1(A)
 \longrightarrow H^1(H)
 \longrightarrow Y_2
 \xrightarrow{\operatorname{tr}}H^2(A)
 \longrightarrow H^2(H).
\]
The first map is an isomorphism: every character of \(H\) factors
through \(U\), hence through \(A\), and \(H\to A\) is surjective.
Kummer theory and Hilbert 90 identify
\(H^2(H,\mathbf F_2)\) with \(\operatorname{Br}(\mathbf Q)[2]\).
This proves \eqref{eq:fr-Y2-transgression}.

For precision, the cohomological exact sequences used here are
the continuous versions of the Hochschild--Serre sequences.
The original cochain construction and transgression are
\cite[Chapter~II, pp.~118--119; Chapter~III, Theorem~2 and the
remark on pp.~129--130]{HochschildSerre1953}.
Continuous cochains and passage through finite quotients are
\cite[II, \S4, Proposition~4.2, pp.~87--88]{MilneCFT}.
The Brauer identification is
\cite[IV, Theorem~3.14, p.~137]{MilneCFT}, combined
with the Kummer exact sequence.

The kernel of \(A\to U\) is
\[
 T=2\widehat{\mathbf Z}.
\]
It has \(H^1(T,\mathbf F_2)=\mathbf F_2 z\), where
\(z(2s)=s_2\bmod2\), and \(H^j(T,\mathbf F_2)=0\) for \(j\ge2\).
Indeed its pro-odd factor has vanishing positive cohomology
by finite averaging, while \(\mathbf Z_2\) has cohomological
dimension one. The latter also follows from the length-one
resolution over \(\mathbf F_2[[\mathbf Z_2]]=\mathbf F_2[[X]]\)
given by multiplication by \(X\).
The action of \(U\) on \(H^1(T,\mathbf F_2)\) is trivial.

Use the continuous section \(u\mapsto(u,\widetilde a(u))\),
where \(\widetilde a(u)\) is the integer \(0\) or \(1\)
representing \(a(u)\). Its extension cocycle, evaluated by \(z\), is
\[
 \frac{\widetilde a(u)+u\widetilde a(v)
                    -\widetilde a(uv)}2\bmod2
 =a(u)a(v)+b(u)a(v).
\]
The class \(a\) factors through the torsion-free
\(\mathbf Z_2\) factor of \(\mathbf Z_2^\times\), so \(a^2=0\).
Thus
\begin{equation}\label{eq:fr-Y2-extension-class}
                         d_2(z)=ab.
\end{equation}

Modulo pro-odd factors, the cyclotomic group is
\[
 U^{(2)}\cong C_2\times\mathbf Z_2
                   \times\prod_{p\ {\rm odd}} C_{2^{k_p}}.
\]
The first two factors have the specified dual characters
\(b\) and \(a\). Its first cohomology has basis
\(b,a,c_p\). Its second cohomology has basis
\begin{equation}\label{eq:fr-Y2-U2-basis}
 \delta_2,\ \delta_p,\quad
 ab,\quad bc_p,\quad ac_p,\quad c_pc_q\ (p<q).
\end{equation}
For a finite cyclic \(2\)-group,
\[
 c_p^2=
 \begin{cases}
  \delta_p,&p\equiv3\pmod4,\\
  0,&p\equiv1\pmod4.
 \end{cases}
\]
These assertions follow from the cyclic resolution and the
Kunneth formula over \(\mathbf F_2\). An explicit check of the
square relation uses the carry cocycle for
\(C_{2^{k+1}}\to C_{2^k}\): when \(k=1\) it is the square of
the degree-one class; when \(k\ge2\) that square is the
coboundary of \(j\mapsto\lfloor j/2\rfloor\bmod2\).
Only finitely many factors occur in any continuous cohomology
class. Thus the displayed basis is a direct-sum basis,
even though the group itself is a product.

The two-row spectral sequence for \(T\to A\to U\), together
with \eqref{eq:fr-Y2-extension-class}, gives
\begin{equation}\label{eq:fr-Y2-H2A}
 0\longrightarrow H^2(U)/\mathbf F_2ab
 \xrightarrow{\operatorname{Inf}_U^A}H^2(A)
 \xrightarrow{\operatorname{edge}}\mathbf F_2a
 \longrightarrow0.
\end{equation}
Here the edge target is the kernel of
\[
 H^1(U)\longrightarrow H^3(U),\qquad x\longmapsto xab.
\]
That kernel is precisely \(\mathbf F_2a\).
The products \(ab^2\) and \(abc_p\) are nonzero independent
classes by their distinct Kunneth multidegrees, whereas
\(a^2b=0\). There are no higher differentials: only the rows
of fibre degrees zero and one are nonzero.
The differential as cup product with the extension class
is \cite[Chapter~III, \S6, Theorems~3--4,
pp.~132--133]{HochschildSerre1953}; its sign is immaterial over
\(\mathbf F_2\).

\subsection{The cubic character is exactly the extra edge direction}

Choose \(\sigma\in H\cap\ker\lambda\) with
\(\rho_2^I(\sigma)=2\). Such an element exists by the
given surjection onto the affine group. Put
\(\alpha=m_3(\sigma)\in\mathbf Z_2\).
The established cubic congruence
\(\rho_{2,2}^I\equiv-2m_3\pmod4\) makes \(\alpha\) odd.

The quotient extension determined by \(\operatorname{tr}(\eta_3)\)
splits over \(T\), using the section
\[
             2s\longmapsto \sigma^s\pmod{\ker\eta_3}.
\]
This is a genuine continuous profinite-power section.
For \(h\in H\) with multiplier \(u=\lambda(h)\), its failure
to be equivariant is measured by
\[
 d(h,s)=h\sigma^s h^{-1}\sigma^{-us}\in K.
\]
The full profinite exponent is \(us\); its dyadic component
is used only when evaluating the cubic coordinate.
The Kummer conjugation law proves \(d(h,s)\in K\).
The additive cubic coordinate on the cyclotomic subgroup
and its weight-three conjugation law give
\[
 \eta_3(d(h,s))
 =\frac{u_2^3-u_2}{8}\,\alpha s_2\bmod2
 =a(u)s_2\bmod2.
\]
The last equality holds because \(u_2\) and \(\alpha\) are odd
and \((u_2^2-1)/8\bmod2=a(u)\).
This is the explicit extension-theoretic description of the edge map.
Consequently
\begin{equation}\label{eq:fr-Y2-eta-edge}
       \operatorname{edge}(\operatorname{tr}(\eta_3))=a.
\end{equation}
In particular the cubic character supplies the extra line in
\eqref{eq:fr-Y2-H2A}; it is not a cyclotomic cup class.
Since its transgression inflates to zero on \(H\), subtracting
its appropriate multiple reduces the computation of \(Y_2\)
to the kernel of the Brauer map on \(H^2(U)/\mathbf F_2ab\).

\subsection{The diagonal classes and their exact local invariants}

For every finite rational prime \(v\), the image of \(\delta_v\)
in \(\operatorname{Br}(\mathbf Q)[2]\) is ramified exactly at
\(\{v,\infty\}\). Here ramification means local invariant \(1/2\).

For an odd prime \(p\), the class \(\delta_p\) is the obstruction
to lifting the maximal \(2\)-power cyclotomic character at \(p\)
from \(C_{2^{k_p}}\) to \(C_{2^{k_p+1}}\).
At any finite place other than \(p\), this character is
unramified and lifts through the procyclic unramified quotient.
At \(p\), local reciprocity identifies its restriction to the
torsion units with a primitive character of the cyclic group
of order \(2^{k_p}\). Such a character cannot lift one step:
an element of order at most \(2^{k_p}\) in
\(C_{2^{k_p+1}}\) reduces to an even element.
Thus the local obstruction is nonzero.
At the real place, complex conjugation maps to the unique
element of order two in \(C_{2^{k_p}}\); it likewise has no
order-two lift. This is the other nonzero local obstruction.

For \(\delta_2=b^2\), the same argument uses the quadratic
character of \(\mathbf Q(i)\). It is unramified at all odd
places; at two its value on the torsion unit \(-1\) cannot
lift from \(C_2\) to \(C_4\); the real obstruction is nonzero.
This proves the assertion for all \(v\).

The precise arithmetic inputs are local reciprocity
\cite[I, Theorem~1.1, p.~20]{MilneCFT}, the local
Brauer invariant \cite[III, Theorem~2.1, p.~103]{MilneCFT},
and the global exact sequence
\cite[VIII, Theorem~4.2, p.~241]{MilneCFT}.
The last theorem identifies \(\operatorname{Br}(\mathbf Q)[2]\)
with the finitely supported tuples of local bits of even total sum.
Therefore the classes ramified at \(\{v,\infty\}\), for finite
primes \(v\), form a basis. In particular the diagonal
subspace spanned by the \(\delta_v\) maps isomorphically onto
\(\operatorname{Br}(\mathbf Q)[2]\).

\subsection{The corrected cross cups}

For a quadratic character of a rational number \(d\), cup
product with the character of \(e\) has local Brauer bit
given by the quadratic Hilbert symbol \((d,e)_v\).
This is the cup-product definition and norm criterion in
\cite[III, \S4, pp.~112--115, especially Theorem~4.4]{MilneCFT}.
We compute only the cases needed for the displayed basis.

First, \((-1,2)\) is trivial globally because
\(2=N_{\mathbf Q(i)/\mathbf Q}(1+i)\).
Thus the zero-Brauer cross cup \(ab\) is precisely the class
already killed in passing from \(U\) to \(A\).

For an odd prime \(p\), the cup \(bc_p\) has no finite
invariant except possibly at \(p\), where its bit is
\(\ell_p(-1)\). There is no invariant at two:
\(p^*\equiv1\pmod4\), so its quadratic extension of
\(\mathbf Q_2\) is split or unramified and every unit is a norm.
Hence the corrected class is
\[
 C_{b,p}=bc_p+\ell_p(-1)\delta_p.
\]

The cup \(ac_p\) can have finite invariants only at \(2\) and \(p\).
At \(p\) its bit is \(\varepsilon_p\).
Its real invariant is zero, so the product formula gives
the same bit at two. Its corrected class is
\[
 C_{2,p}=ac_p+\varepsilon_p(\delta_2+\delta_p).
\]

Finally, for distinct odd primes \(p,q\), the cup \(c_pc_q\)
has finite bits
\[
 \ell_p(q^*)\quad\text{at }p,\qquad
 \ell_q(p^*)\quad\text{at }q,
\]
and zero at every other finite place. At two both parameters
are odd and congruent to one modulo four, so the same
split-or-unramified norm observation applies.
The corrected class is
\[
 C_{p,q}=c_pc_q+\ell_p(q^*)\delta_p+\ell_q(p^*)\delta_q.
\]
The real bit is automatically cancelled as well, by the
global sum formula.

For completeness, the odd-prime local calculations just used
are elementary. A quadratic extension defined by a unit is
split or unramified, so a second unit is a norm. A quadratic
extension defined by a parameter of odd valuation is tamely
ramified, and its unit norms have square residues. Thus a
uniformizer against a unit gives the corresponding Legendre
symbol. There are no unlisted finite places.

Adding the uniquely determined diagonal correction to each
cross cup is a triangular change of the basis
\eqref{eq:fr-Y2-U2-basis}. Since the diagonal subspace maps
isomorphically onto the Brauer group, it follows that
\[
 \ker\!\left(H^2(U)\longrightarrow\operatorname{Br}(\mathbf Q)[2]\right)
 =
 \mathbf F_2ab
 \oplus\bigoplus_p\mathbf F_2 C_{b,p}
 \oplus\bigoplus_p\mathbf F_2 C_{2,p}
 \oplus\bigoplus_{p<q}\mathbf F_2 C_{p,q}.
\]
Quotienting by \(ab\), and adjoining the extra cubic line
from \eqref{eq:fr-Y2-eta-edge}, proves the asserted completeness
and independence once the known Kummer coordinates are identified.

\subsection{Identification of the odd-prime Kummer axes}

Let \(\epsilon_p\in H^1(U,\mathbf F_2)\) denote the character
of \(\mathbf Q(\sqrt p)\), with the positive rational prime
as radicand. Then
\[
 \epsilon_p=
 \begin{cases}
 c_p,&p\equiv1\pmod4,\\
 b+c_p,&p\equiv3\pmod4.
 \end{cases}
\]
The mod-four Kummer cocycle of \(p\), whose mod-two value is
\(\epsilon_p\), determines a central quadratic extension over \(A\).
Using its integer-valued parity section, exactly the same carry
calculation as above gives
\[
 \operatorname{tr}(\psi_p)
 =\operatorname{Inf}_U^A(\epsilon_p^2+b\epsilon_p).
\]
If \(p\equiv1\pmod4\), this is \(\operatorname{Inf}(bc_p)\).
If \(p\equiv3\pmod4\), it is
\(\operatorname{Inf}(bc_p+c_p^2)
 =\operatorname{Inf}(bc_p+\delta_p)\).
Both cases give
\[
                  \operatorname{tr}(\psi_p)
                         =\operatorname{Inf}_U^A(C_{b,p}).
\]
This proves every assertion in
Theorem~\ref{thm:fr-Y2-complete}.

\subsection{A completely explicit target outside the known span}

The character \(\kappa_{5,29}\) has a small radical realization.
Put
\[
 \alpha=7+2\sqrt5,\qquad
 L=\mathbf Q(\sqrt5,\sqrt{29},\sqrt\alpha).
\]
Its norm from \(\mathbf Q(\sqrt5)\) is \(29\).
Let \(F=\mathbf Q(\sqrt5)\) and \(E=F(\sqrt{29})\).
Neither \(\alpha\) nor \(\alpha/29\) is a square in \(F\),
because its norm is respectively \(29\) or \(1/29\),
neither a square in \(\mathbf Q\).
An element of \(F\) becomes a square in \(F(\sqrt{29})\)
only if it or its quotient by \(29\) was already a square in \(F\):
expand \((x+y\sqrt{29})^2\) and use \(2xy=0\).
Thus \([L:E]=2\), so \([L:\mathbf Q]=8\).

Write \(c=\sqrt\alpha\), \(d=\sqrt{29}\).
The transformations
\[
 r:\quad \sqrt5\mapsto-\sqrt5,\ d\mapsto d,\ c\mapsto d/c,
 \qquad
 s:\quad \sqrt5\mapsto\sqrt5,\ d\mapsto-d,\ c\mapsto c
\]
are involutions. Their commutator fixes \(E\) and sends
\(c\) to \(-c\). Hence \(L/\mathbf Q\) is dihedral of order
eight, with central quadratic subgroup \(\operatorname{Gal}(L/E)\).

Choose \(u^{(5)},u^{(29)}\in U\) supported at the indicated
coefficient primes, with value \(2\) at that prime and \(1\)
elsewhere. The number \(2\) is a nonsquare modulo both
\(5\) and \(29\). Their dyadic components are one, so
\[
 (u^{(5)},0),\qquad (u^{(29)},0)
\]
are commuting elements of \(A\).
Choose arbitrary lifts \(h_5,h_{29}\in H\).
Then
\[
                         k=[h_5,h_{29}]\in K.
\]
On \(L\), these lifts have the two independent images of
\(r\) and \(s\) modulo the centre. Multiplying a lift by a
central element does not change their commutator.
Consequently \(k\) sends \(c\) to \(-c\).

This proves directly that \(L\not\subset K_2\).
Since \(E\subset\mathbf Q^{\rm ab}\subset K_2\) and
\([L:E]=2\), one has \(L\cap K_2=E\).
The quadratic extension \(K_2L/K_2\) therefore gives a
nonzero \(H\)-invariant character \(\chi_L\in Y_2\),
with \(\chi_L(k)=1\).

Every known \(\psi_p\) vanishes on \(k\).
Indeed both lifts have dyadic multiplier one, so each
dyadic Kummer cocycle is a homomorphism on the subgroup
they generate and is zero on their commutator.
Likewise the cubic coordinate is a homomorphism on
\(\ker\lambda_2\): the dyadic degree-two term vanishes
when \(\lambda_2=1\), and the degree-three term is additive
under the GT product. Hence \(m_3(k)=0\) and
\(\eta_3(k)=0\).
Thus
\[
 \chi_L\notin
 \mathbf F_2\eta_3+\bigoplus_p\mathbf F_2\psi_p.
\]

Finally \((5,29)\) is split, because \(29=N_{F/\mathbf Q}(\alpha)\).
The two coordinate involutions in the dihedral presentation
have square one and nontrivial commutator. Its central
extension class is therefore the cross cup \(c_5c_{29}\).
It follows that \(\chi_L=\kappa_{5,29}\) in the normalization
of \eqref{eq:fr-Y2-kpq}. The norm construction of such dihedral
fields is classical; see \cite[Lemma~5.1 and Corollary~5.2,
author pp.~9--10]{StevenhagenRedei2022}.

The argument proves no lift of \(\chi_L\), or of any other
new target basis vector, to the full-GT invariant character
space. That remains a separate extension problem.

